\documentclass[11pt]{article}
\usepackage[T1]{fontenc}
\usepackage{lmodern}
\usepackage{amsmath,amssymb,amsthm}
\usepackage{graphicx}
\usepackage[margin=1in]{geometry}
\usepackage{microtype}
\usepackage[colorlinks=true,linkcolor=blue,citecolor=blue,urlcolor=blue]{hyperref}

\newtheorem{theorem}{Theorem}
\newtheorem{lemma}{Lemma}
\newcommand{\E}{\mathbb E}
\newcommand{\PP}{\mathbb P}
\newcommand{\Om}{\Omega}

\title{The Neumann stochastic square in the planar CH/AC homotopy}
\author{Candidate solution packet for arXiv:2204.09545}
\date{13 August 2026}

\begin{document}
\maketitle

\begin{center}
\fbox{\parbox{0.91\linewidth}{\textbf{Status.} Candidate full answer for
smooth bounded planar domains, likely valid, awaiting specialist review.}}
\end{center}

\section{Source problem}

Bl\"omker and T\"olle isolate the following missing estimate in Section 6 of
\cite{BT}, on page 18 of the source PDF.

\begin{center}
\includegraphics[width=\linewidth]{figures/open_problem_crop.png}
\end{center}

We prove the requested convergence for every connected bounded
$C^\infty$ domain in $\mathbb R^2$.  The proof also quantifies the exact
noise scaling and limiting constant.

\section{Statement}

Let $L=-\Delta_N$ be the nonnegative Neumann Laplacian on $L^2(\Om)$.
Write
\[
 Le_j=\mu_j e_j,\qquad
 e_0=|\Om|^{-1/2},\quad \mu_0=0,
\]
and work first on the mean-zero subspace, so indices below start at $1$.
For $0<\varepsilon<1/2$, put
\[
 a_\varepsilon=1-\varepsilon,
 \qquad B_\varepsilon=L(a_\varepsilon+\varepsilon L),
 \qquad
 \lambda_j^\varepsilon=\mu_j(a_\varepsilon+\varepsilon\mu_j).
\]
If $W_0$ is cylindrical white noise on the mean-zero subspace, define the
zero-initial stochastic convolution
\[
 Z_\varepsilon(t)
 =\sigma_\varepsilon\int_0^t
   e^{-(t-s)B_\varepsilon}\,dW_0(s).
\]
For fixed $\varepsilon$, the fourth-order regularization makes the ordinary
square $Z_\varepsilon^2$ well-defined.

We use the Bessel-potential norm
$\|f\|_{H^{-1}}=\|(I+L)^{-1/2}f\|_2$.

\begin{theorem}\label{thm:main}
Suppose
\[
 \sigma_\varepsilon^2\log(1/\varepsilon)\longrightarrow\kappa
 \quad\text{for some }\kappa\in[0,\infty).
\]
Then for every $T<\infty$, $1\leq p<\infty$, and
$q\geq\max\{2,p\}$,
\[
 \left(\E\left[
 \left\|Z_\varepsilon^2-\frac{\kappa}{8\pi}\right\|_
 {L^p(0,T;H^{-1}(\Om))}^{q}
 \right]\right)^{1/q}\longrightarrow0.
\]
More precisely, with constants depending on $\Om,T,p,q$ but not on
$\varepsilon$,
\begin{equation}\label{eq:rate}
 \left(\E\left[
 \left\|Z_\varepsilon^2-\frac{\kappa}{8\pi}\right\|_
 {L^pH^{-1}}^{q}
 \right]\right)^{1/q}
 \leq C_{p,q,T}
 \left(\left|\sigma_\varepsilon^2\log(1/\varepsilon)-\kappa\right|
       +\sigma_\varepsilon^2\right).
\end{equation}
In particular the convergence holds in probability in every finite
$L^p(0,T;H^{-1})$.
\end{theorem}

The source suppresses the constant Fourier mode on the torus.  The same
convention is used here.  Restoring the one-dimensional Neumann kernel adds
terms that are $O(\sigma_\varepsilon)$ in the relevant Gaussian norms and
does not alter the conclusion.

\section{Proof intuition}

Although products of Neumann eigenfunctions have no Fourier convolution law,
the linear operator is still a function of the Neumann Laplacian.  Its
covariance is therefore an integral of the ordinary Neumann heat kernel.  On
the diagonal, the interior heat kernel contributes the universal
$(8\pi)^{-1}\log(1/\varepsilon)$ divergence.  The reflected boundary term is
large only in a shrinking collar and leaves an $L^2$-bounded remainder.

After subtracting the mean, Wick's formula turns the $H^{-1}$ variance of the
square into the integral of three kernels: one Bessel Green kernel and two
covariance kernels.  Each has only a logarithmic diagonal singularity, and a
logarithmic cube is integrable in two dimensions.  The remaining factor is
$\sigma_\varepsilon^4$, so the random part vanishes.  This replaces the
torus-specific Fourier cancellation by a domain-independent heat-kernel
argument.

\section{Two heat-kernel estimates}

Let $p_s(x,y)$ be the Neumann heat kernel and
$p_s^0(x,y)=p_s(x,y)-|\Om|^{-1}$.  We use two standard consequences of the
Gaussian estimate and the not-feeling-the-boundary estimate
\cite{LiStrohmaier}.  If $d(x)=\operatorname{dist}(x,\partial\Om)$, then,
for $0<s\leq1$,
\begin{align}
 |p_s(x,y)|&\leq Cs^{-1}e^{-|x-y|^2/(Cs)},\label{eq:gauss}\\
 \left|p_s(x,x)-(4\pi s)^{-1}\right|
 &\leq Cs^{-1}e^{-d(x)^2/(Cs)}+Cs^{-1/2}.\label{eq:nfb}
\end{align}
The second term is the harmless lower-order parametrix remainder.  For
$s\geq1$, $p_s^0$ decays exponentially by the spectral gap.

Define the unit-amplitude stationary covariance kernel
\begin{equation}\label{eq:Qdef}
 Q_\varepsilon(x,y)
 =\sum_{j\geq1}\frac{e_j(x)e_j(y)}{2\lambda_j^\varepsilon}.
\end{equation}

\begin{lemma}\label{lem:heat}
Uniformly for $0<\varepsilon<1/2$,
\begin{align}
 Q_\varepsilon(x,x)
 &=\frac{1}{8\pi}\log(1/\varepsilon)+R_\varepsilon(x),
 &\sup_\varepsilon\|R_\varepsilon\|_{L^2(\Om)}&<\infty,
 \label{eq:diag}\\
 |Q_\varepsilon(x,y)|
 &\leq C\bigl(1+|\log|x-y||\bigr),&&x\neq y.
 \label{eq:offdiag}
\end{align}
The kernel $G_1$ of $(I+L)^{-1}$ obeys the same off-diagonal logarithmic
bound and is nonnegative.
\end{lemma}

\begin{proof}
For $\mu>0$, partial fractions and the Laplace transform give
\[
 \frac1{\mu(a_\varepsilon+\varepsilon\mu)}
 =\frac1{a_\varepsilon}
   \int_0^\infty(1-e^{-a_\varepsilon s/\varepsilon})e^{-\mu s}\,ds.
\]
Consequently,
\begin{equation}\label{eq:heat-rep}
 Q_\varepsilon(x,y)
 =\frac1{2a_\varepsilon}\int_0^\infty
 (1-e^{-a_\varepsilon s/\varepsilon})p_s^0(x,y)\,ds.
\end{equation}
Estimate \eqref{eq:gauss}, followed by the change of variables
$u=|x-y|^2/s$, proves \eqref{eq:offdiag}.  The same heat representation,
without the cutoff factor, proves the assertion for $G_1$.

On the diagonal, substitute $(4\pi s)^{-1}$ in the part of
\eqref{eq:heat-rep} with $0<s<1$.  Since $a_\varepsilon\in[1/2,1]$,
\[
 \frac1{8\pi a_\varepsilon}\int_0^1
 (1-e^{-a_\varepsilon s/\varepsilon})\frac{ds}{s}
 =\frac1{8\pi}\log(1/\varepsilon)+O(1).
\]
The large-time and constant-mode terms are uniformly bounded.  By
\eqref{eq:nfb}, the absolute value of the remaining boundary term is bounded
by a constant plus
\[
 H_\varepsilon(d(x))
 :=\int_0^1\min\{1,s/\varepsilon\}\,
 s^{-1}e^{-d(x)^2/(Cs)}\,ds.
\]
For $d\leq\sqrt\varepsilon$, this is at most $C\log(1/\varepsilon)$;
for $d>\sqrt\varepsilon$, it is at most $C(1+|\log d|)$.  Smooth boundary
normal coordinates therefore give
\[
 \sup_{\varepsilon<1/2}\int_\Om H_\varepsilon(d(x))^2\,dx<\infty,
\]
because $\sqrt\varepsilon\log^2(1/\varepsilon)$ is bounded and
$\int_0^1|\log r|^2dr<\infty$.  This proves \eqref{eq:diag}.
\end{proof}

\section{The deterministic mean}

Set $m_\varepsilon(t,x)=\E Z_\varepsilon(t,x)^2$.  Spectral calculus gives
\begin{equation}\label{eq:mean}
 m_\varepsilon(t,x)
 =\sigma_\varepsilon^2 Q_\varepsilon(x,x)
 -D_\varepsilon(t,x),
\quad
 D_\varepsilon(t,x)
 =\sigma_\varepsilon^2\sum_{j\geq1}
 \frac{e^{-2t\lambda_j^\varepsilon}}{2\lambda_j^\varepsilon}e_j(x)^2.
\end{equation}
Lemma~\ref{lem:heat} shows
\begin{equation}\label{eq:stationary-mean}
 \left\|\sigma_\varepsilon^2Q_\varepsilon(\cdot,\cdot)|_{x=y}
 -\frac\kappa{8\pi}\right\|_{H^{-1}}
 \leq C\left(
 |\sigma_\varepsilon^2\log(1/\varepsilon)-\kappa|
 +\sigma_\varepsilon^2\right).
\end{equation}

Since $\lambda_j^\varepsilon\geq\mu_j/2$,
\[
 0\leq \frac{D_\varepsilon(t,x)}{\sigma_\varepsilon^2}
 \leq\sum_{j\geq1}\frac{e^{-t\mu_j}}{\mu_j}e_j(x)^2
 =\int_t^\infty p_s^0(x,x)\,ds
 \leq C_T(1+|\log t|),
\]
where the final bound follows from \eqref{eq:gauss} and the spectral gap.
Thus, for every finite $p$,
\begin{equation}\label{eq:mean-bound}
 \left\|m_\varepsilon-\frac\kappa{8\pi}\right\|_{L^p(0,T;H^{-1})}
 \leq C_{p,T}\left(
 |\sigma_\varepsilon^2\log(1/\varepsilon)-\kappa|
 +\sigma_\varepsilon^2\right).
\end{equation}

\section{The centered square}

Let $Y_\varepsilon(t)=Z_\varepsilon(t)^2-m_\varepsilon(t)$.  First truncate
the Neumann expansion at a finite index.  Wick's formula gives
\begin{align}
 \E\|Y_\varepsilon(t)\|_{H^{-1}}^2
 &=2\sum_{i,j\geq1}c_i(t)c_j(t)
 \left\langle e_ie_j,(I+L)^{-1}(e_ie_j)\right\rangle,
 \label{eq:wicksum}
\end{align}
where
\[
 c_j(t)=\frac{\sigma_\varepsilon^2
 (1-e^{-2t\lambda_j^\varepsilon})}{2\lambda_j^\varepsilon}
 \leq\frac{\sigma_\varepsilon^2}{2\lambda_j^\varepsilon}.
\]
Every summand in \eqref{eq:wicksum} is nonnegative.  Hence monotone passage
from finite truncations and the kernel identity yield
\begin{align*}
 \E\|Y_\varepsilon(t)\|_{H^{-1}}^2
 &\leq2\sigma_\varepsilon^4
 \iint_{\Om^2}G_1(x,y)Q_\varepsilon(x,y)^2\,dx\,dy\\
 &\leq C\sigma_\varepsilon^4.
\end{align*}
The last integral is uniformly finite by Lemma~\ref{lem:heat}, since
$(1+|\log|x-y||)^3$ is locally integrable in two dimensions.

The random variable $Y_\varepsilon(t)$ takes values in the second homogeneous
Wiener chaos.  Vector-valued Gaussian hypercontractivity therefore gives,
for every $q\geq2$,
\[
 \left(\E\|Y_\varepsilon(t)\|_{H^{-1}}^q\right)^{1/q}
 \leq C_q\sigma_\varepsilon^2.
\]
For $q\geq p$, Minkowski's integral inequality now implies
\begin{equation}\label{eq:centered}
 \left(\E\|Y_\varepsilon\|_{L^p(0,T;H^{-1})}^q\right)^{1/q}
 \leq C_{p,q,T}\sigma_\varepsilon^2.
\end{equation}
Combining \eqref{eq:mean-bound} and \eqref{eq:centered} proves
Theorem~\ref{thm:main}.

\section{Scope, verification, and novelty}

The theorem supplies precisely the domain-dependent estimate isolated in the
source paper, for arbitrary smooth shape rather than a square or torus.  It
also clarifies that the nontrivial regime is expressed invariantly as
$\sigma_\varepsilon^2\log(1/\varepsilon)\to\kappa$; whether a paper denotes
noise amplitude or variance by $\sigma_\varepsilon$ can otherwise obscure a
square root.

The smooth-boundary hypothesis is used only in the boundary-collar form of
the heat-kernel estimate.  We make no claim here for rough domains or colored
noise operators that fail to commute with $L$.  A numerical companion checks
the coefficient $1/(8\pi)$ from Neumann eigenvalues on the unit square; it is
not part of the proof.

A bounded novelty search through 13 August 2026 covered the run registry,
local full sources, the exact open phrase, arXiv:2204.09545, and web searches
combining CH/AC homotopy, Neumann boundary, general domain, stochastic
convolution, and Wick square.  No later explicit resolution was found.  This
is not an exhaustive priority search.

\paragraph{Human-review recommendation.}
Verify first Lemma~\ref{lem:heat}, especially the uniform $L^2$ treatment of
the boundary collar.  Next verify the nonnegative spectral expansion
\eqref{eq:wicksum} and its kernel identity.  These are the only substantive
points; the remaining bounds are direct.

\begin{thebibliography}{9}
\bibitem{BT}
D. Bl\"omker and J. M. T\"olle,
\emph{Singular limits for stochastic equations},
Stochastics and Dynamics (2023), DOI: 10.1142/S0219493723500405;
arXiv:2204.09545.

\bibitem{LiStrohmaier}
L. Li and A. Strohmaier,
\emph{Heat kernel estimates for general boundary problems},
J. Spectr. Theory 6 (2016), 903--919;
arXiv:1604.00784.
\end{thebibliography}

\end{document}
